%
En billøfter på et verksted virker som en hydraulisk jekk. Det lille stempelet har diameter \SI{2.0}{\centi\metre}, og det store stempelet har diameter \SI{20}{\centi\metre}. En bil med masse \SI{1500}{\kilo\gram} står på det store stempelet. Vi ser bort fra massen til stemplene og antar at begge stemplene er i samme høyde.
\begin{enumerate}[a)]
    \item Hvor mange ganger større er arealet av det store stempelet enn arealet av det lille? (Arealet av en sirkel er $A = \pi r^2$.)
    \item Hvor stor kraft må virke på det lille stempelet for å holde bilen oppe?
    \item Det lille stempelet presses \SI{10}{\centi\metre} ned. Hvor mye løftes bilen? \textit{Hint: Væsken som presses ut av den lille sylinderen, må gå inn i den store.}
    \item Regn ut arbeidet som gjøres på det lille stempelet, og arbeidet som gjøres på bilen. Hva ser du? Får vi noe «gratis» med en hydraulisk jekk?
\end{enumerate}
%
